\chapter{Conclusion and Outlook\label{cha:chapter7}}

The final chapter summarizes the thesis, reviews the main empirical findings, outlines practical contributions, and provides an outlook on future research.

\section{Summary of Contributions\label{sec:summary}}

This research examines the effect of prompt framing on code generation and mathematical reasoning by large language models (LLMs). It conducts an extensive full factorial experiment that examines several factors, including request tone, role framing, and prompt wrapper design. The study covers six local Ollama models, two API models, two types of tasks, 50 base tasks per task type, 18 prompt designs, and a total of 14,400 experiments.

The main conclusion is that model ability dominates correctness. At the same time, how the instructions are worded or designed can still have an impact on the results and behavior of a model, although the effects of the prompt are less significant than those of the model and vary from one task to another or based on the model itself. The core mixed-effects logistic analysis reveals that structured prompts improve API Python performance, with the XML format showing the largest gain, while JSON decreases local math performance. Moreover, after clustered-robustness and dominant-model sensitivity checks, no drastic changes in the direction of the effects were observed.

The practical consequence is that prompt design should be treated as an empirical engineering parameter in software systems based on LLMs. Effective prompt strategies cannot be derived from folklore. They should be tested against the actual variables that influence them: the specific model, task types, parser, evaluator, time limits, and token price.

For practitioners, the best default option is not a fixed phrasing, but rather a process: practitioners should start with a neutral explicit baseline, define the expected outcome, measure the correctness and behavior together, and only then proceed to revise the prompt. Structured wrappers are promising for some coding workflows, mainly because they reduce irrelevant text and improve parser compatibility, though they should not be blindly transferred to final-answer math tasks.

\section{Data Availability and Replication Package\label{sec:replication}}

The aggregated result dataset is publicly available on Kaggle \cite{promptframingkaggle2026}. Prompt templates, task metadata, evaluation scripts, analysis notebooks, and aggregated result files are maintained in the project repository \cite{promptframingrepo2026}.

Key thesis replication artifacts include:
\begin{itemize}
\tightlist
\item
 \href{https://github.com/tycoi2005/prompt-framing-experiment/blob/main/02ScalingPhrase/results_analysis/prompt_framing_experiment_results.jsonl}{Aggregated result JSONL}
\item
 \href{https://github.com/tycoi2005/prompt-framing-experiment/blob/main/02ScalingPhrase/results_analysis/notebooks/smallmodels_analysis.executed.ipynb}{Local-model analysis notebook}
\item
 \href{https://github.com/tycoi2005/prompt-framing-experiment/blob/main/02ScalingPhrase/results_analysis/notebooks/apimodels_analysis.executed.ipynb}{API-model analysis notebook}
\item
 \href{https://github.com/tycoi2005/prompt-framing-experiment/blob/main/02ScalingPhrase/results_analysis/notebooks/smallmodels_without_nemotron_analysis.executed.ipynb}{No-Nemotron sensitivity notebook}
\item
 \href{https://github.com/tycoi2005/prompt-framing-experiment/blob/main/10Paperv3/scripts/run_mixed_effects_analysis.py}{Mixed-effects analysis script}
\end{itemize}

\section{Statements and Declarations}

\subsection{Funding Statement}
This work received no specific grant from any funding agency, commercial, or not-for-profit sectors.

\subsection{Competing Interests}
The author declares no competing interests.

\subsection{AI Tools Declaration}
AI tools were used for editing support, code assistance, and drafting assistance. The experimental outputs were generated by the evaluated models and were not manually altered for the reported analyses.

\subsection{Ethics Statement}
This study does not involve human participants. It analyzes model outputs on coding and mathematical reasoning tasks.

\section{Outlook and Future Work\label{sec:outlook}}

Further research ought to investigate larger and more diverse code benchmarks, other API models, mixed-effects models with richer random-effect structures, raw-response-preserving math audits, and user/workflow studies with developers applying prompt policies in realistic coding workflows. Multi-turn agent settings are especially important since repair loops, tool feedback, and constrained decoding may change the role of prompt framing. Recent work on multimodal LLMs \cite{yin2024multimodal} and their applications in various domains \cite{zhang2025cybersecurity} suggests that prompt sensitivity research should extend beyond text-only interactions. The general aim is to take prompt engineering from folklore to reproducible system design, treating prompts as versioned, tested parts of an LLM system rather than isolated incantations, which makes the design auditable and easier to improve.
